% GENERATED FILE. Edit canonical lessons, then run npm run generate:books.
% canonical-source-hash: fnv1a64:6e9ba35785af24e0
% canonical-lessons: SA-C499-pitrsvasa, SA-C499-matrsvasa, SA-C499-bhratrjah, SA-C499-bhagineyah, SA-C499-dauhitrah

\chapter{Aunts and Nephews}
\label{ch:sa-aunts-and-nephews}
\hlchaptermodality{499}

\begin{chapteropening}
\textbf{By the end of this chapter:} \emph{I can name more members of the family.}

Complete the last lesson of chapter 499: I can name more members of the family.
\end{chapteropening}

\section[\texorpdfstring{\sk{पितृष्वसा} (pitṛṣvasā)}{पितृष्वसा (pitṛṣvasā)}]{\sk{पितृष्वसा} (pitṛṣvasā) — a father's sister}

\label{lesson:SA-C499-pitrsvasa}



\begin{quote}
Before the new one: say the Sanskrit for youngest, then the Sanskrit for twins.
\end{quote}

\subsection*{You'll want to know: \sk{पितृष्वसा}}
\textbf{\sk{पितृष्वसा}} — \emph{pitṛṣvasā} — \textquotedblleft{}a father's sister\textquotedblright{}.

\begin{quote}
\sk{मम} \sk{पितृष्वसा} \sk{ग्रामे} \sk{वसति।}
\end{quote}

\begin{quote}
\emph{mama pitṛṣvasā grāme vasati.}
\end{quote}

\begin{quote}
My father's sister lives in the village.
\end{quote}

A feminine noun.

\begin{cousinweb}
\textbf{\sk{पितृष्वसा}} (\emph{pitṛṣvasā}) is built from \emph{pitṛ}, \textquotedblleft{}father\textquotedblright{} (\textbf{\sk{पिता}}), and \emph{svasṛ}, \textquotedblleft{}sister\textquotedblright{} (MW); CDIAL 8177 lists it beside \emph{pituḥṣvasṛ}, the form behind Bengali \emph{pisi} (CDIAL 8177).
\end{cousinweb}

\begin{grammarlens}[title={What you've built}]
One more word for describing what you lived and felt.
\end{grammarlens}

\subsection*{Your turn}
\begin{itemize}
\raggedright
  \item \emph{Say it:} \emph{pitṛṣvasā}
  \item \emph{Say it:} \emph{pitṛṣvasā}, once more
  \item \emph{Say it:} \emph{mama pitṛṣvasā grāme vasati.} — My father's sister lives in the village.
  \item \emph{Recall:} say the Sanskrit for youngest, then the Sanskrit for twins, then say \emph{pitṛṣvasā} again
\end{itemize}

\begin{tcolorbox}[breakable,colback=teal!4,colframe=teal!35!black,title={Before you move on}]
What does \sk{पितृष्वसा} mean? (\textquotedblleft{}A father's sister\textquotedblright{}.) Say it once more.
\end{tcolorbox}
\section[\texorpdfstring{\sk{मातृष्वसा} (mātṛṣvasā)}{मातृष्वसा (mātṛṣvasā)}]{\sk{मातृष्वसा} (mātṛṣvasā) — a mother's sister}

\label{lesson:SA-C499-matrsvasa}



\begin{quote}
Before the new one: say the Sanskrit for twins, then the Sanskrit for a father's sister.
\end{quote}

\subsection*{You'll want to know: \sk{मातृष्वसा}}
\textbf{\sk{मातृष्वसा}} — \emph{mātṛṣvasā} — \textquotedblleft{}a mother's sister\textquotedblright{}.

\begin{quote}
\sk{मातृष्वसा} \sk{मह्यं} \sk{मोदकान्} \sk{ददाति।}
\end{quote}

\begin{quote}
\emph{mātṛṣvasā mahyaṃ modakān dadāti.}
\end{quote}

\begin{quote}
My mother's sister gives me sweets.
\end{quote}

A feminine noun.

\begin{cousinweb}
\textbf{\sk{मातृष्वसा}} (\emph{mātṛṣvasā}) is built from \emph{mātṛ}, \textquotedblleft{}mother\textquotedblright{} (\textbf{\sk{माता}}), and \emph{svasṛ}, \textquotedblleft{}sister\textquotedblright{} (MW); CDIAL 10001 lists it beside \emph{mātuḥṣvasṛ}, which lives on as Hindi \emph{mausī} (CDIAL 10001).
\end{cousinweb}

\begin{grammarlens}[title={What you've built}]
One more word for describing what you lived and felt.
\end{grammarlens}

\subsection*{Your turn}
\begin{itemize}
\raggedright
  \item \emph{Say it:} \emph{mātṛṣvasā}
  \item \emph{Say it:} \emph{mātṛṣvasā}, once more
  \item \emph{Say it:} \emph{mātṛṣvasā mahyaṃ modakān dadāti.} — My mother's sister gives me sweets.
  \item \emph{Recall:} say the Sanskrit for twins, then the Sanskrit for a father's sister, then say \emph{mātṛṣvasā} again
\end{itemize}

\begin{tcolorbox}[breakable,colback=teal!4,colframe=teal!35!black,title={Before you move on}]
What does \sk{मातृष्वसा} mean? (\textquotedblleft{}A mother's sister\textquotedblright{}.) Say it once more.
\end{tcolorbox}
\section[\texorpdfstring{\sk{भ्रातृजः} (bhrātṛjaḥ)}{भ्रातृजः (bhrātṛjaḥ)}]{\sk{भ्रातृजः} (bhrātṛjaḥ) — a brother's son}

\label{lesson:SA-C499-bhratrjah}



\begin{quote}
Before the new one: say the Sanskrit for a father's sister, then the Sanskrit for a mother's sister.
\end{quote}

\subsection*{You'll want to know: \sk{भ्रातृजः}}
\textbf{\sk{भ्रातृजः}} — \emph{bhrātṛjaḥ} — \textquotedblleft{}a brother's son\textquotedblright{}.

\begin{quote}
\sk{मम} \sk{भ्रातृजः} \sk{विद्यालये} \sk{पठति।}
\end{quote}

\begin{quote}
\emph{mama bhrātṛjaḥ vidyālaye paṭhati.}
\end{quote}

\begin{quote}
My brother's son studies at school.
\end{quote}

A masculine noun.

\begin{cousinweb}
\textbf{\sk{भ्रातृजः}} (\emph{bhrātṛjaḥ}) is built from \emph{bhrātṛ}, \textquotedblleft{}brother\textquotedblright{} (\textbf{\sk{भ्राता}}), and \emph{ja}, \textquotedblleft{}born\textquotedblright{}, from the root √jan; MW gives \textquotedblleft{}a brother's son\textquotedblright{} (MW).
\end{cousinweb}

\begin{grammarlens}[title={What you've built}]
One more word for describing what you lived and felt.
\end{grammarlens}

\subsection*{Your turn}
\begin{itemize}
\raggedright
  \item \emph{Say it:} \emph{bhrātṛjaḥ}
  \item \emph{Say it:} \emph{bhrātṛjaḥ}, once more
  \item \emph{Say it:} \emph{mama bhrātṛjaḥ vidyālaye paṭhati.} — My brother's son studies at school.
  \item \emph{Recall:} say the Sanskrit for a father's sister, then the Sanskrit for a mother's sister, then say \emph{bhrātṛjaḥ} again
\end{itemize}

\begin{tcolorbox}[breakable,colback=teal!4,colframe=teal!35!black,title={Before you move on}]
What does \sk{भ्रातृजः} mean? (\textquotedblleft{}A brother's son\textquotedblright{}.) Say it once more.
\end{tcolorbox}
\section[\texorpdfstring{\sk{भागिनेयः} (bhāgineyaḥ)}{भागिनेयः (bhāgineyaḥ)}]{\sk{भागिनेयः} (bhāgineyaḥ) — a sister's son}

\label{lesson:SA-C499-bhagineyah}



\begin{quote}
Before the new one: say the Sanskrit for a mother's sister, then the Sanskrit for a brother's son.
\end{quote}

\subsection*{You'll want to know: \sk{भागिनेयः}}
\textbf{\sk{भागिनेयः}} — \emph{bhāgineyaḥ} — \textquotedblleft{}a sister's son\textquotedblright{}.

\begin{quote}
\sk{मातुलः} \sk{भागिनेयाय} \sk{उपहारं} \sk{ददाति।}
\end{quote}

\begin{quote}
\emph{mātulaḥ bhāgineyāya upahāraṃ dadāti.}
\end{quote}

\begin{quote}
The uncle gives his sister's son a present.
\end{quote}

A masculine noun.

\begin{cousinweb}
\textbf{\sk{भागिनेयः}} (\emph{bhāgineyaḥ}) is built on \textbf{\sk{भगिनी}}, \textquotedblleft{}a sister\textquotedblright{}, by lengthening its first vowel; MW gives \textquotedblleft{}a sister's son\textquotedblright{} (MW).
\end{cousinweb}

\begin{grammarlens}[title={What you've built}]
One more word for describing what you lived and felt.
\end{grammarlens}

\subsection*{Your turn}
\begin{itemize}
\raggedright
  \item \emph{Say it:} \emph{bhāgineyaḥ}
  \item \emph{Say it:} \emph{bhāgineyaḥ}, once more
  \item \emph{Say it:} \emph{mātulaḥ bhāgineyāya upahāraṃ dadāti.} — The uncle gives his sister's son a present.
  \item \emph{Recall:} say the Sanskrit for a mother's sister, then the Sanskrit for a brother's son, then say \emph{bhāgineyaḥ} again
\end{itemize}

\begin{tcolorbox}[breakable,colback=teal!4,colframe=teal!35!black,title={Before you move on}]
What does \sk{भागिनेयः} mean? (\textquotedblleft{}A sister's son\textquotedblright{}.) Say it once more.
\end{tcolorbox}
\section[\texorpdfstring{\sk{दौहित्रः} (dauhitraḥ)}{दौहित्रः (dauhitraḥ)}]{\sk{दौहित्रः} (dauhitraḥ) — a daughter's son}

\label{lesson:SA-C499-dauhitrah}



\begin{quote}
Before the new one: say the Sanskrit for a brother's son, then the Sanskrit for a sister's son.
\end{quote}

\subsection*{You'll want to know: \sk{दौहित्रः}}
\textbf{\sk{दौहित्रः}} — \emph{dauhitraḥ} — \textquotedblleft{}a daughter's son\textquotedblright{}.

\begin{quote}
\sk{मातामहः} \sk{दौहित्रेण} \sk{सह} \sk{क्रीडति।}
\end{quote}

\begin{quote}
\emph{mātāmahaḥ dauhitreṇa saha krīḍati.}
\end{quote}

\begin{quote}
Grandfather plays with his daughter's son.
\end{quote}

A masculine noun.

\begin{cousinweb}
\textbf{\sk{दौहित्रः}} (\emph{dauhitraḥ}) is built on \emph{duhitṛ}, \textquotedblleft{}a daughter\textquotedblright{} (\textbf{\sk{दुहिता}}), by lengthening its first vowel; MW gives \textquotedblleft{}a daughter's son\textquotedblright{} (MW).
\end{cousinweb}

\begin{grammarlens}[title={What you've built}]
One more word for describing what you lived and felt.
\end{grammarlens}

\subsection*{Your turn}
\begin{itemize}
\raggedright
  \item \emph{Say it:} \emph{dauhitraḥ}
  \item \emph{Say it:} \emph{dauhitraḥ}, once more
  \item \emph{Say it:} \emph{mātāmahaḥ dauhitreṇa saha krīḍati.} — Grandfather plays with his daughter's son.
  \item \emph{Recall:} say the Sanskrit for a brother's son, then the Sanskrit for a sister's son, then say \emph{dauhitraḥ} again
\end{itemize}

\begin{tcolorbox}[breakable,colback=teal!4,colframe=teal!35!black,title={Before you move on}]
What does \sk{दौहित्रः} mean? (\textquotedblleft{}A daughter's son\textquotedblright{}.) Say it once more.
\end{tcolorbox}
